Unmanned aerial vehicles (UAVs) have become very popular tools across military, commercial, and research sectors.\cite{drones} Accurate aerodynamic analysis is critical for optimizing their performance, particularly when evaluating the impact of external payloads on drag and efficiency. Computational Fluid Dynamics (CFD) offers a cost-effective alternative to wind tunnel testing for aerodynamic predictions. However, CFD results must be properly validated against experimental data before they can be reliably applied to new design configurations.

In aerospace engineering, commercial CFD programmes such as ANSYS Fluent dominate due to their convenient user interface and comprehensive support. OpenFOAM, an open-source CFD toolbox, provides a compelling alternative with zero costs and complete source code access. Despite its capabilities, OpenFOAM remains underutilized in aerospace applications compared to commercial software, partly due to difficult learning process and concerns about validation. 

This thesis addresses this gap by validating OpenFOAM against both wind tunnel experimental data and ANSYS Fluent predictions for a small-scale delta wing aircraft. The validation approach uses published wind tunnel measurements of a delta wing configuration with dimensions comparable to the target drone developed by CAFA Tech. By comparing aerodynamic coefficients from OpenFOAM simulations against both experimental results and ANSYS Fluent computations, the accuracy of the open-source solver can be quantitatively assessed. Delta wing configurations are particularly suitable for validation due to their complex vortical flow structures that challenge CFD accuracy.

Once the OpenFOAM simulation framework is validated, the same computational setup will be applied to analyze the CAFA Tech fixed-wing drone in two configurations: the baseline clean aircraft and the configuration with two rocket-shaped external stores mounted beneath the fuselage. The primary objective is to evaluate a drag increase by the external stores. This information is essential for CAFA Tech's high speed drone development, because external payloads reduce aerodynamic efficiency. 
